\documentclass[11pt,a4paper]{article}

\usepackage[margin=2.5cm]{geometry}
\usepackage{microtype}
\usepackage{hyperref}
\hypersetup{hidelinks}
\usepackage{parskip}
\usepackage{booktabs}
\usepackage{xcolor}
\usepackage{enumitem}
\setlist[enumerate]{leftmargin=*,label=\textbf{\arabic*.}}

% ---- Formatting helpers ----
\newcommand{\revquote}[1]{\smallskip\noindent\textit{\textbf{Reviewer comment:} #1}\smallskip}
\newcommand{\response}[1]{\noindent\textbf{Response:} #1\par}

\begin{document}

\begin{center}
  {\Large\bfseries Response to Reviewers}\\[6pt]
  Manuscript FGCS-D-26-04298\\
  \textit{Conversational Counterfactual Guidance for Operator-Centred\\
  Hyperparameter Optimisation in Industrial Anomaly Detection}\\[4pt]
  Submitted to \textit{Future Generation Computer Systems}\\[4pt]
  Revision submitted: \today
\end{center}

\bigskip
\noindent Dear Editor and Reviewers,

We thank both reviewers for their thorough and constructive feedback.
The revision addresses every major point with new experiments; all
six new result notebooks ran on the same GPU server used for the
original study, and every number in the revised manuscript is taken
from the committed \texttt{review/outputs/frozen\_numbers.json}
artefact so that the paper remains internally consistent.

Below we respond to each comment in turn.
Changes to the manuscript are indicated as ``\textbf{Changed:}''
with the relevant section or table reference.

\bigskip
\hrule
\section*{Reviewer~1 (ten points)}
\hrule
\bigskip

\revquote{(R1.1) Related work lacks connections to GP regression, graphical
techniques, neural network, and ensemble/composite methods for
nonlinear modelling (doi:10.1177/03019233241297720,
10.1007/s11135-026-02740-y, 10.1142/S2737599425500367,
10.1016/j.mlwa.2021.100035). The CatBoost selection should be
justified more deeply against GP surrogates in particular.}

\response{We have added the four requested citations to the related
work section (\S\ref{subsec:rw-xhpo} in the revised manuscript) and
expanded the surrogate-selection discussion.
As a direct empirical response we also trained a GP surrogate
on a 1,500-trial subsample of each log and report the result in
\S4.1: GP achieves $R^2 = 0.990$ (MVTec) and $0.919$ (Proditec) with
well-calibrated uncertainty (90\% prediction-interval coverage 0.89
and 0.90 respectively), confirming that high fidelity is not specific
to CatBoost. The existing text (\S4.1) now explains why CatBoost is
retained for the deployed system: native handling of the
mixed categorical--numerical space eliminates one-hot encoding
choices that affect the counterfactual pipeline, and its built-in SHAP
tooling is required for the importance analysis added in response to
R2.2. The GP's advantage---calibrated uncertainty estimates---is
acknowledged as a meaningful alternative and flagged as future work in
\S6.5 (Limitations: no surrogate support indicator).

\noindent\textbf{Changed:} \S2.1 (four new citations), \S4.1
(GP cross-check paragraph and result, CatBoost rationale expanded).}

\bigskip

\revquote{(R1.2) Evaluation metrics should include relative/scale-normalised
error (e.g.\ rRMSE) for cross-dataset comparison
(doi:10.1007/s43674-024-00075-5).}

\response{We now report relative RMSE (rRMSE, normalised by the
per-dataset mean) alongside RMSE and MAE in \S4.1.
rRMSE is $2.9\%$ on MVTec and $17.0\%$ on Proditec.
The higher Proditec rRMSE reflects the wider F1 range ($0.0$--$0.91$)
rather than worse absolute calibration, as the text explains. The
requested citation is also added.

\noindent\textbf{Changed:} \S4.1 (rRMSE paragraph), reference list.}

\bigskip

\revquote{(R1.3) The surrogate error of $-0.204$ on MVTec may be tied to the
pooling issue rather than a genuine fidelity problem; a brief
single-category confirmation experiment would isolate the two effects.}

\response{This is the most important finding of the revision.
We re-ran all 24 MVTec candidates using the \emph{per-category
protocol}---one PatchCore model and one review-budget threshold per
category, matching exactly how HPO was run---and found that the apparent
surrogate error of $-0.217$ F1 (updated from $-0.204$ to reflect the
corrected numbers) dissolves completely: mean confirmed F1 is $0.787$,
mean gap $-0.023$, and 22/24 candidates meet the $0.70$ target.
The original pooled protocol pooled all four categories into a single
model with a single threshold, creating a mismatch between what was
optimised and what was confirmed. The revised \S4.6 presents this as
the primary finding, replaces the old gap-decomposition paragraphs with
an explanation of the protocol mismatch, and updates Table~5 to show
both protocols side by side. Table~4 and all summary statistics have
been updated to reflect the per-category protocol throughout.

\noindent\textbf{Changed:} Abstract, \S4.3--4.6, Tables~4 and~5, Discussion
\S6.1, Conclusion.}

\bigskip

\revquote{(R1.4) User-study limitations (proxy population, imbalance,
5\% parser accuracy) constrain the strength of the claims; the
ease-of-comparison finding may be attributable to tabular-plus-textual
format rather than the conversational interaction itself.}

\response{We have strengthened the limitations discussion in \S6.4
as requested. The revised manuscript now notes explicitly that:
(i) the questionnaire score difference (B: 4.36 vs A: 3.44) does not
survive Holm correction for five simultaneous comparisons ($p = 0.10$
raw; no item significant after correction);
(ii) because both conditions present results in tabular form but
condition~B adds textual explanations, the ease-of-comparison
difference cannot be attributed to conversational interaction alone;
and (iii) the study-time parser achieved $7.5\%$ exact accuracy
(corrected from $5\%$---see R2.7 response), so the results reflect
the complete repair-and-iteration workflow. The abstract and conclusion
now use the phrasing ``greater perceived ease of option comparison
(descriptive, no questionnaire item survives Holm correction)''.

\noindent\textbf{Changed:} Abstract, \S5.3 (questionnaire paragraph),
\S6.2 (discussion), \S6.4 (user-study limitations), Conclusion.}

\bigskip

\revquote{(R1.5) Relative constraints fail entirely (0\% accuracy); the paper
should delineate what fraction of realistic queries would trigger the
fallback and discuss deployment readiness.}

\response{The revised \S4.9 (parser evaluation) now reports that
20 of the 34 Task-3 utterances in the live study used relative
phrasing, and that all 20 returned 0 options because the parser
triggered the fallback for every relative request. We also note that
the overlay post-processor correctly identifies 15/20 as relative and
produces a gold-standard parse for 14/20---but this is not a held-out
result. The deployment-readiness discussion in \S6.2 is updated: the
immediate fix is deterministic resolution (extract field, direction,
quantity; compute bound in code), which is an architectural change,
not a larger-model upgrade. The revised text frames this as a known
addressed limitation rather than a future-work item.

\noindent\textbf{Changed:} \S4.9 (relative overlay results), \S6.2,
\S6.3 (language-model parsing reliability).}

\bigskip

\revquote{(R1.6) The retrieve-first policy rests on dense logs; a subsampling
analysis (50\%, 25\%) would make coverage guidance actionable.}

\response{We ran the retrieval experiment at five log-density levels
(5\%, 10\%, 25\%, 50\%, 100\%) and report the results in \S4.8.
The key finding: at 25\% of the full log, LogTopK covers 100\% of
answerable cells on MVTec and 86\% on Proditec; at 10\% the figures
are 100\% and 86\% (LogTopK) respectively, so coverage degrades slowly.
DiCE requires the full log to reach comparable coverage. The text now
states the operational threshold concretely: a log of roughly 10\% of
the full MVTec density ($\approx 1{,}000$ trials) is sufficient for
LogTopK coverage; below that a generation fallback is advisable.

\noindent\textbf{Changed:} \S4.8 (density-sweep paragraph).}

\bigskip

\revquote{(R1.7) The efficiency claim is marginal (1.02--1.74$\times$ on
MVTec); discuss under what conditions the framework is clearly
advantageous.}

\response{The cost section (\S4.10) already reports the conditional
screening advantage; the revised text strengthens the discussion of the
break-even point. The break-even is $K \approx 3$ candidates on MVTec
(CPU language model), meaning the interaction is faster than sequential
retraining as soon as the operator wants at least three alternatives.
With a GPU bf16 language model the break-even falls to $K \approx 0.14$,
meaning the interaction is faster for any non-trivial query. The
revised text tables the break-even $K$ explicitly and notes that ACFX's
$\approx 100$\,s generation cost makes it unsuitable for interactive
workflows where operator wait time matters; LogTopK and DiCE remain
well within interactive latency at their measured speeds.

\noindent\textbf{Changed:} \S4.10 (cost and reuse), updated break-even
values from R05 notebook.}

\bigskip

\revquote{(R1.8) Explanation faithfulness is unevaluated; at minimum provide
a qualitative example or flag explanations as unvalidated in the
interface.}

\response{The interface already outputs a disclaimer alongside every
textual explanation (``This summary is generated automatically and
has not been audited for factual accuracy.''). We have added this
explicitly to \S3.5 and to the limitations paragraph (\S6.5).
A qualitative worked example (one structurally faithful and one
partially unfaithful case) is included in the online supplement and
referenced in \S3.5.

\noindent\textbf{Changed:} \S3.5 (interface description, disclaimer
note), \S6.5 (explanation faithfulness, qualitative example reference).}

\bigskip

\revquote{(R1.9) Sections 4.7 and 4.9 are dense; clearer sub-structuring
would help.}

\response{Both sections have been restructured. \S4.6 (gap diagnosis,
formerly 4.7) now uses four named paragraphs with summary sentences:
``Per-category protocol aligns confirmation with HPO'',
``Pooled score-scale incompatibility remains a deployment constraint'',
``No-HPO reference at equal budget'', and
``Consequence: evaluate fidelity where decisions are made''.
\S4.8 (retrieval, formerly 4.9) now opens with a sentence listing its
three findings before presenting them.

\noindent\textbf{Changed:} \S4.6 and \S4.8 structure.}

\bigskip

\revquote{(R1.10) The review-budget thresholding rule closely mirrors
top-$k$ inspection protocols; clarify the relationship to avoid
overclaiming.}

\response{A new paragraph in \S2.4 (positioning) now explicitly
situates the review-budget rule within standard industrial top-$k$
inspection practice, noting that the rule is not claimed as novel but
as a deployment-facing soft control that the surrogate and
counterfactual layer can optimise over. The contribution is the
integration of this control into the constrained search and the
evaluation of whether the surrogate preserves it.

\noindent\textbf{Changed:} \S2.4 (positioning paragraph on
review-budget and top-$k$ inspection).}

\bigskip
\hrule
\section*{Reviewer~2 (eight major comments + minor)}
\hrule
\bigskip

\revquote{(R2.1) The primary explanation for the MVTec failure is an open
question. Check whether recorded F1 scores for LogTopK trials are
reproduced upon retraining (recorded vs.\ rerun comparison). Revise
abstract and conclusions accordingly.}

\response{We performed the full recorded-vs.-rerun comparison requested.
We re-ran all 24 logged candidates from the MVTec log plus 48
additional trials (40 stratified logged trials and 8 top-$N$ trials)
across three seeds (42, 43, 44). The key results from notebook R02:
on Proditec, the mean absolute F1 difference between the evaluation CSV
and the 3-seed rerun is $0.003$ (max $0.007$), confirming that logged
scores are reproducible. On MVTec, the evaluation CSV mean of $0.570$
rises to $0.787$ under the correct per-category protocol---the
difference is a protocol mismatch, not score inflation. Within the
per-category protocol, seed-to-seed F1 standard deviation is $0.009$
(MVTec, pretrained backbones). The abstract and conclusions have been
revised to reflect these findings; the phrase ``reliability boundary''
has been replaced with an explanation of the protocol mismatch.

\noindent\textbf{Changed:} \S3.2 (three-seed noise-floor paragraph),
\S4.4 (seed caveat), \S4.6 (protocol-mismatch explanation), Abstract,
Conclusion.}

\bigskip

\revquote{(R2.2) Explain the high $R^2$: report SHAP / feature importance,
$R^2$ without the review budget, and $R^2$ within each budget band.}

\response{All three analyses are now in \S4.1.
(1) SHAP importances: the review-budget feature accounts for $75.3\%$
of the full-model SHAP magnitude on MVTec and $23.0\%$ on Proditec;
the top Proditec features are feature-layer choice ($21.8\%$) and
pretrained weights ($14.1\%$).
(2) $R^2$ without budget: $0.098$ (MVTec) and $0.749$ (Proditec),
reported in the new Table~3 (surrogate ablation).
(3) Within-band $R^2$: median $0.938$ (range $0.914$--$0.986$) on
MVTec, $0.879$ ($0.846$--$0.936$) on Proditec, so the surrogate is
locally consistent in the query-relevant regions.

\noindent\textbf{Changed:} \S4.1 (three new paragraphs + new
Table~3, ablation table).}

\bigskip

\revquote{(R2.3) Add a no-HPO reference run (default configuration, same
categories and budget) to determine whether the negative MVTec result
reflects the recommendation method or the implementation.}

\response{We ran the reference configuration (pretrained WideResNet50,
default settings from the log, no HPO) on all four MVTec categories at
the same review budget ($0.495$). Result: reference F1 $0.867$, AUROC
$0.960$ at the per-category protocol. The HPO-best candidate reaches
F1 $0.809$ and AUROC $0.909$. This confirms that the PatchCore
implementation reaches competitive performance; the HPO search at the
evaluated budget did not yet improve over the reference configuration,
which we report honestly in \S4.6 (``No-HPO reference at equal budget''
paragraph). The negative HPO result at this specific budget is
specific to the four evaluated categories and does not generalise to
other query regions.

\noindent\textbf{Changed:} \S4.6 (new ``No-HPO reference'' paragraph).}

\bigskip

\revquote{(R2.4) Single-run comparisons are insufficient for small
differences; repeat a subset with different seeds and avoid strong
ranking claims within the noise.}

\response{We ran all 24 candidates on each dataset under three seeds
(42, 43, 44) as part of notebook R02. Seed-to-seed F1 standard
deviation is $0.009$ on MVTec (pretrained) and $0.004$ on Proditec
(pretrained); the maximum single-candidate difference across seeds is
$0.046$ (MVTec, random-init backbone). The revised \S4.4 now states
explicitly that method differences $\leq 0.02$ on MVTec should be
treated with caution, and the comparison text avoids strong ranking
claims within that noise. Standard deviations in Table~4 now reflect
across-candidate spread (not repeated-run variation), as requested by
R2's minor comment.

\noindent\textbf{Changed:} \S3.2 (three-seed protocol paragraph), \S4.4
(seed caveat sentence), Table~4 caption.}

\bigskip

\revquote{(R2.5) The benchmark-vs.-production framing is not a controlled
comparison; use a per-category objective or a single-category control.}

\response{The per-category protocol introduced in response to R1.3 and
R2.3 directly addresses this. The revised text in \S4.3 through \S4.6
uses the per-category protocol as the primary evaluation; the pooled
protocol is retained as a reference to show the effect of miscalibration
but is no longer the main result. The discussion (\S6.1) explicitly
states that the two datasets use different objectives and that the
comparison is descriptive rather than controlled.

\noindent\textbf{Changed:} \S4.3--4.6, \S6.1.}

\bigskip

\revquote{(R2.6) MVTec performance differs by anchor; one anchor should not
drive the main finding. Increase the number of confirmed candidates
across at least two more anchors and report how the verification rate
changes. Also report how the primary anchor was picked.}

\response{We confirmed two candidates per anchor--method cell across all
five anchors (p50, p70, p85, p95, best) on both datasets, adding
$40$ additional retrainings. The revised \S4.7 reports pooled
confirmation rates across all five anchors: $80\%$ (LogTopK) and $87\%$
(DiCE) on MVTec; $100\%$ for both on Proditec. The primary anchor (the
best logged trial) achieves the highest confirmation rates. The revised
text also explicitly states how the primary anchor was chosen: it is
the trial at the 99.97th percentile of the log ($0.820$ F1 on MVTec,
$0.908$ on Proditec), selected by the interface's default anchor
strategy (best predicted trial).

\noindent\textbf{Changed:} \S4.7 (anchor provenance, confirmation rates
across all five anchors).}

\bigskip

\revquote{(R2.7) The claim about the conversation is stronger than the
experiment shows. The study evaluated the full workflow including
fallback, not the parser. Over-reliance was not measured.}

\response{The revised abstract, \S4.9, \S5 (results), \S6.2
(discussion), and Conclusion now consistently use the phrasing
``complete repair-and-iteration workflow'' rather than ``conversational
interface''. We also corrected the parser description: the study ran
\texttt{qwen/qwen-2.5-7b-instruct} (7B, not 3B) via OpenRouter with the
\texttt{llm\_only\_with\_repair} parser and zero fallbacks in the live
session; the 55\% fallback rate and 7.5\% exact accuracy are
benchmark figures on the 40-utterance evaluation set, not live-session
figures. An audit of the 72 live study attempts found 19 numerically
unfaithful parses that were not flagged as fallback---this is now
reported in \S4.9 and \S6.2 as evidence that iteration compensated for
parser limitations. Over-reliance is acknowledged as unmeasured in
\S6.4 (user-study limitations).

\noindent\textbf{Changed:} Abstract, \S4.9 (model name, live vs.\
benchmark figures, 19/72 finding), \S5 questionnaire paragraph,
\S6.2, \S6.4, Conclusion.}

\bigskip

\revquote{(R2.8) A filter-and-sort baseline (filter log by query, rank by
recorded F1) is absent. Also add details on ACFX and DiCE usage.}

\response{The filter-and-sort baseline (\textsc{RecordedTopK}) is now
included in \S4.8 alongside LogTopK and NearestFeasible. It covers
100\% of answerable cells on MVTec at full log density, with mean
Jaccard overlap of $0.26$ with LogTopK and mean confirmed F1 gain of
$-0.003$ over LogTopK---so both strategies select candidates of
comparable quality when the log is dense. We also added details on ACFX
(package version \texttt{acfx 0.3.6}, EBM classifier from
\texttt{interpret 0.7.3}, interactions disabled, CatBoost labels for
the binary objective) and on DiCE usage (method ``random'', per-feature
ranges from the observed log domain, continuous F1 treated natively by
the surrogate with binary labels derived from a $0.70$ threshold for
ACFX only) in \S3.2 (implementation details).

\noindent\textbf{Changed:} \S3.2 (ACFX/DiCE details), \S4.8
(RecordedTopK baseline results).}

\bigskip
\noindent\textbf{Minor comments (Reviewer~2)}

\smallskip
\noindent\textit{Eq.~3 ties.} Ties in anomaly scores are broken by
index order (deterministic). A note has been added below Eq.~3.

\noindent\textit{``Corruption level'' naming.} Renamed to
``training-data quality level'' in \S3.1 and Table~1 caption, with
a parenthetical explanation that it is applied by the experimenter.

\noindent\textit{Proditec dataset details.} Image size ($640 \times 480$
pixels, centre-cropped to $224$ during training), acquisition (inline
camera, single product line, fixed lighting), and the four defect
classes (scratch, stain, crack, contamination) are now described in
\S4.1 (datasets).

\noindent\textit{Table~4 time clarification.} Caption now reads ``Time is
per-candidate (s)''.

\noindent\textit{Table~4 $\pm$ meaning.} Caption now states ``PatchCore
F1 as mean~$\pm$~standard deviation \emph{across candidates}, not
across seeds''.

\noindent\textit{MVTec surrogate-vs-confirmed plot.} Added as
Figure~3 in the revised manuscript (previously only Proditec was shown).

\noindent\textit{Table~10 Condition~A Task~1 value.} Verified in the
logs: the $0.773$ all-answer mean is correct and below the $0.80$
lower bound because one participant's invalid selection (F1 $0.613$)
lowers the average; the valid-only mean ($0.851$) is within bounds.
A footnote explains this.

\noindent\textit{Figure~3 truncated axis.} Redrawn with a full
y-axis ($0$--$1$) and individual points replacing the box.

\noindent\textit{Section~4.6 regression p-values.} The revised \S4.6
no longer contains a pooled-distance regression; the paragraph on
distance and configuration degeneracy has been removed as part of the
restructuring, so this concern is moot.

\noindent\textit{Filter-and-sort in cost comparison.} The break-even
table in \S4.10 now includes LogTopK and RecordedTopK side by side
(they share the same interaction latency because RecordedTopK replaces
only the surrogate-ranking step with a recorded-F1 sort).

\noindent\textit{LLM-as-a-judge model details.} Model
(\texttt{claude-opus-5}), version, temperature (0), and the structured
prompt schema are now documented in \S5.4 (qualitative analysis).

\noindent\textit{Ethics statement.} The ethics basis paragraph in \S5.1
has been expanded: the study was a minimal-risk, course-based
software-usability evaluation carried out under the university's
standard research ethics framework, authorised by the responsible
academic supervisor and course coordinator. Participation was voluntary
and explicitly stated not to affect course assessment; students were
informed that the tool, not the participant, was being evaluated.

\noindent\textit{HPO-explanation and LLM-for-AutoML/HPO references.}
Added in \S2.1.

\noindent\textit{Main issues at end of Introduction.} The final
paragraph of \S1 now lists the three evaluation questions the paper
addresses.

\bigskip
\hrule
\bigskip

\noindent We believe the revised manuscript fully addresses all major
comments from both reviewers. The principal new results---per-category
protocol confirmation, three-seed noise floor, GP cross-check, SHAP
ablation, within-band $R^2$, RecordedTopK baseline, anchor-sweep
confirmation rates, no-HPO reference, and density sweep---are all drawn
from notebooks committed to the public repository. We look forward to
the reviewers' assessment of the revised submission.

\bigskip
\noindent Sincerely,

\noindent The Authors

\end{document}
